\documentclass[11pt]{article}
\usepackage[margin=0.8in]{geometry}
\usepackage[T1]{fontenc}
\usepackage[utf8]{inputenc}
\usepackage{booktabs,longtable,hyperref,fancyvrb,parskip}
\usepackage[expansion=false]{microtype}

\hypersetup{
  colorlinks=true,
  linkcolor=black,
  urlcolor=black,
  citecolor=black,
  pdftitle={LeafOS llama.cpp Vulkan Provider Spine},
  pdfauthor={LeafOS}
}

\title{\textbf{LeafOS llama.cpp Vulkan Provider Spine}\\[4pt]\large Runtime Boundary and Operator Checks}
\author{LeafOS}
\date{}

\newcommand{\code}[1]{\texttt{#1}}

\begin{document}
\maketitle

\section*{Purpose}

This document records the intended local inference boundary for LeafOS:
\code{llama.cpp} is the model provider, and Vulkan is the preferred local GPU
backend. The provider is a long-lived \code{llama-server} process. It is not the
node transport and it is not the filesystem authority.

\section*{Control Plane}

\begin{Verbatim}[fontsize=\small, frame=single]
Human task
  -> task file
  -> CPU graph and policy checks
  -> provider call to llama.cpp HTTP endpoint
  -> model proposal
  -> parser
  -> patch gate / verification gate
  -> accepted artifact or rejection
  -> journal / report
\end{Verbatim}

\section*{Provider Startup}

On Windows, use the PowerShell provider launcher:

\begin{Verbatim}[fontsize=\small, frame=single]
cd C:\path\to\LeafOS\PowerShell-Version
.\serve-llamacpp.ps1 --llama-dir C:\llama-vulkan --backend vulkan --model fable
\end{Verbatim}

Use \code{--check-only} before a real run:

\begin{Verbatim}[fontsize=\small, frame=single]
.\serve-llamacpp.ps1 --llama-dir C:\llama-vulkan --check-only --backend vulkan
\end{Verbatim}

\section*{Provider Environment}

Set these in the shell that will run LeafOS commands:

\begin{Verbatim}[fontsize=\small, frame=single]
$env:LEAF_PROVIDER_MODE = 'llamacpp'
$env:LEAF_BRAIN_PROVIDER = 'llamacpp'
$env:LEAF_CODER_PROVIDER = 'llamacpp'
$env:LEAF_LLAMACPP_URL = 'http://127.0.0.1:8080'
\end{Verbatim}

\section*{Diagnostics}

\begin{longtable}{@{}p{0.34\linewidth}p{0.56\linewidth}@{}}
\toprule
Check & Meaning \\
\midrule
\code{vulkaninfo --summary} & Confirms the OS Vulkan stack sees a GPU. \\
\code{llama-server --list-devices} & Confirms the selected llama.cpp build exposes accelerator devices. \\
\code{leafctl provider test llamacpp} & Confirms the HTTP provider is reachable and returns model text. \\
\code{nvidia-smi} & Confirms the NVIDIA driver and live GPU activity. \\
\bottomrule
\end{longtable}

\section*{Common Failure}

\begin{Verbatim}[fontsize=\small, frame=single]
Observed:
  vulkaninfo sees NVIDIA GeForce RTX 4070
  llama-server --list-devices prints only:
    Available devices:

Meaning:
  The selected llama.cpp package is CPU-only.

Fix:
  Use a Windows x64 Vulkan llama.cpp build.
  Pass its directory with --llama-dir or set LEAF_LLAMACPP_DIR.
\end{Verbatim}

\section*{Runtime Doctrine}

\begin{Verbatim}[fontsize=\small, frame=single]
The GPU proposes.
The CPU disposes.

Model output is untrusted text.
The parser turns it into a typed artifact.
The patch gate decides whether it can touch files.
The journal records what happened.
\end{Verbatim}

\section*{Promotion Gate}

The provider spine is ready for \code{0.4.0-alpha.1} evidence when all of these
are true:

\begin{enumerate}
  \item \code{serve-llamacpp.ps1 --check-only --backend vulkan} sees at least one device.
  \item \code{llama-server} stays alive as a persistent HTTP endpoint.
  \item \code{leafctl provider test llamacpp} passes all checks.
  \item Brain and coder lanes can route to \code{llamacpp} independently.
  \item A generated patch is validated before any apply operation.
\end{enumerate}

\end{document}
